\section{Multithreaded Reference Orbit Acceleration}
\label{subsec:ref-orbit-mt}

The \code{MT3} reference-orbit path accelerates the high-cost MPIR arithmetic
of the single-threaded authoritative orbit (\cref{subsec:ref-orbit-st}) by
decomposing each iteration across multiple CPU threads.  The design targets
independent, high-latency big-integer operations within one reference
recurrence, while minimizing synchronization overhead.  Iterations themselves
remain sequential, but squaring the next imaginary component can begin before
the current real update finishes. Two
recurring implementation patterns characterize this approach.

The \code{MT5} entry point is currently disabled and falls back to \code{MT3};
it is not a distinct active backend.

\paragraph{Asynchronous squaring.}
During authoritative orbit evaluation, two worker threads compute the squared
terms \(x_n^2\) and \(y_n^2\) concurrently.  In parallel, the main thread
evaluates the cross term in
\[
  y_{n+1} = 2 x_n y_n + c_y,
\]
performs low-precision periodicity checks, and manages operand reuse and
serialization.  Once the squared terms are returned, the real component update
is completed as
\begin{equation}
  x_{n+1} = x_n^2 - y_n^2 + c_x.
  \label{eq:mt-real-update}
\end{equation}
The products in a batch all read $x_n,y_n$ before those values are overwritten.
The workers return their products before the main thread publishes the new
state.  For $x_n=3/4$, $y_n=1/2$, and $c=0$, the batch returns $9/16$,
$1/4$, and $2x_ny_n=3/4$, giving $z_{n+1}=5/16+3i/4$.
Using the new real component in the imaginary product would instead give
$2(5/16)(1/2)=5/16$, a different recurrence.  Parallelism therefore requires
stable old operands and a join before their buffers can be reused.

\paragraph{MT3 imaginary-square lookahead.}
The main thread forms $y_{n+1}$ from the cross product before waiting for
$x_n^2$. Once the current $y_n^2$ has been consumed into separate storage,
the Y worker starts $y_{n+1}^2$ while the caller waits for $x_n^2$ and finishes
\eqref{eq:mt-real-update}. A local flag records that the next Y-square is
already in flight. The orbit dependency is therefore finer than a whole-iteration
barrier: $y_{n+1}^2$ needs neither $x_n^2$ nor $x_{n+1}$. Stable square inputs and
separate completed-result storage protect that overlap.

The schedule evaluates the same polynomial as the single-threaded factored
path, but its separately rounded squares need not give bitwise-identical
results at finite precision.

\paragraph{Lock-free handoff with prefetch.}
Inter-thread communication is implemented using a minimal lock-free mailbox
(\code{ThreadPtrs<T>}) consisting of atomic \code{In} and \code{Out}
pointers.  The protocol proceeds as follows:
\begin{enumerate}
  \item The producer fills stable operand buffers and publishes a work pointer
        by storing it into \code{In}.
  \item The worker spins until it atomically claims the pointer, prefetches the
        work descriptor, executes the required arithmetic, and then
        publishes the same pointer via \code{Out}.
  \item The producer spins until it retrieves the completed pointer from
        \code{Out}, after which the computed results are consumed and the
        buffers may be reused.  A published pointer is a handoff, not permission
        for both threads to mutate the same operands concurrently.
\end{enumerate}
To reduce cache-miss latency, \code{PrefetchHighPrec} touches MPIR metadata
and limb data using a fixed 64-byte stride.  In authoritative \code{MT3},
the main thread prefetches completed square results; the explicit square-worker
operand-prefetch calls are inactive.  The reuse perturbation workers also
prefetch their operands.  Prefetching can hide some memory latency, but its
benefit depends on the precision and cache traffic.

\subsection{Precision-dependent single-thread versus multi-thread trade-offs}
\label{subsec:ref-orbit-mt-tradeoff}

The relative effectiveness of single-threaded and multithreaded reference-orbit
evaluation depends strongly on the active precision, which in practice is
closely correlated with zoom level and MPIR limb count.  Consequently, the
implementation supports both execution modes and selects between them
automatically when \code{PerturbationAlg::Auto} is enabled.

At lower precision (small limb counts), overall performance is often dominated
by fixed overheads such as synchronization, atomic mailbox operations, thread
spinning, and cache traffic.  In this regime, a single-threaded MPIR execution
path benefits from tight instruction locality, reduced fencing, and predictable
cache reuse, and may outperform a parallelized approach despite executing all
arithmetic serially.

As precision increases, the cost structure shifts toward throughput-dominated
big-integer arithmetic.  Large limb counts cause MPIR squaring and multiplication
to dominate total runtime, amortizing synchronization overhead and making
concurrent execution increasingly effective.  Overlapping independent
operations, such as the simultaneous computation of \(x_n^2\) and \(y_n^2\),
yields substantial gains, while explicit limb prefetching further mitigates
memory latency.  Beyond a hardware- and workload-dependent threshold,
multithreaded reference-orbit computation can outperform the
single-threaded path in practice, with speedups increasing as precision grows.

The current \code{Auto} implementation first estimates precision in bits,
excluding reuse padding. At $16{,}384$ bits and above it selects
\code{MTPeriodicity9Karatsuba}, without a hardware-concurrency test.
Below that precision threshold it selects \code{STPeriodicity} below $10^{150}$,
\code{MTPeriodicity3} from $10^{150}$ through values below $10^{2000}$, and
\code{MTPeriodicity3PerturbMTHighMTMed3} at $10^{2000}$ and beyond.  These are
fixed code thresholds, not a hardware benchmark or CUDA capability test, and
\code{Auto} never selects the GPU backend.

\paragraph{A simple cost model.}
Let $C_{\mathrm{mul}}$ be one full-precision multiplication's cost,
$C_{\mathrm{sq}}$ a square's cost, $C_{\mathrm{low}}$ the main thread's
reduced-precision work, $C_{\mathrm{join}}$ the handoff and waiting overhead,
and $C_{\mathrm{assemble}}$ the final additions and publication.  Ignoring
memory contention and optional reuse work, the critical paths are approximately
\begin{align}
T_{\mathrm{ST}} &\approx 2C_{\mathrm{mul}}+C_{\mathrm{low}}
                  +C_{\mathrm{assemble}}, \\
T_{\mathrm{MT3}} &\approx
  \max(C_{\mathrm{sq}},C_{\mathrm{mul}}+C_{\mathrm{low}})
  +C_{\mathrm{join}}+C_{\mathrm{assemble}}.
\end{align}
The maximum appears because the three product tasks overlap; the final
assembly waits for the slowest.  As an illustrative calculation, costs of
$12$, $10$, $1$, $2$, and $1$ time units respectively give
$T_{\mathrm{ST}}\approx26$ and $T_{\mathrm{MT3}}\approx16$.
If multiplication and squaring each cost only $0.05$ units but joining still
costs $0.2$, with $0.01$ each for reduced work and assembly, the corresponding
costs are $0.12$ and $0.27$.  These are examples of the model, not measured
benchmarks.  More cores cannot eliminate the join, and shared memory bandwidth
can make the parallel products slower than isolated ones.

Overall, single-threaded and multithreaded reference-orbit evaluation are
treated as complementary strategies rather than a strict hierarchy: fewer limbs
favor single-thread execution, while increasing precision increasingly favors
multithreaded decomposition, with \code{Auto} providing an adaptive runtime
selection mechanism.

\subsection{MT9 exact split-product backend}
\label{subsec:ref-orbit-mt9}

The nine-thread Karatsuba reference design is inspired by
Fractadyne~\cite{Fractadyne}, specifically its \code{ParSplit} and
\code{split\_mandel\_walk} implementation in \code{backend\_rug.rs}.
Each of the recurrence's three products is decomposed into three integer
products of approximately half the original operand width. Eight persistent
workers and the caller provide the nine arithmetic roles.

\paragraph{Precision distinction from Fractadyne.}
Both implementations reconstruct the split integer products exactly.
Fractadyne additionally uses a derivative-driven adaptive working-precision
schedule: when the schedule lowers precision, its recurrence products,
additions, and subtractions round at that reduced precision. FractalShark's
MT9 backend keeps recurrence addition and subtraction at the full configured
MPIR precision throughout the orbit. Its derivative remains a reduced-precision
periodicity estimate; it does not lower the arithmetic precision. In particular,
``half-width'' describes integer multiplication operands, not a reduction of
the recurrence's configured precision. The comparison refers to the inspected
Fractadyne revision identified in the bibliography.

\paragraph{Exact algebra and final truncation.}
Let $\beta=2^w$, where $w$ is the number of usable bits in an MPIR limb, and
let $m$ be the low-part limb count. For the full stored nonnegative input
mantissas, including guard limbs, define
\[
 X=X_H\beta^m+X_L,\qquad Y=Y_H\beta^m+Y_L.
\]
With exact integer products
$A=X_HY_H$, $B=X_LY_L$, and $C=(X_H+X_L)(Y_H+Y_L)$,
\begin{equation}
 XY=A\beta^{2m}+(C-A-B)\beta^m+B.
 \label{eq:mt9-karatsuba}
\end{equation}
The middle term is $X_HY_L+X_LY_H\ge0$; no truncated subtraction or
unaccounted borrow occurs. Squares use the same identity with $X=Y$.
An extra sum limb accommodates $X_H+X_L$, and unequal high/low widths are
handled by zero padding. The complete integer product is assembled before
any output limbs are discarded.

For a nonzero exact floating-point product with normalized limb exponent $e$,
retaining $q$ limbs toward zero gives
\begin{equation}
 |\widehat{xy}-xy|<\beta^{e-q}.
 \label{eq:mt9-local-truncation}
\end{equation}
The destination retains up to \code{\_mp\_prec+1} limbs. Karatsuba introduces
no error before that final truncation. The bound concerns products of the
stored inputs; it does not bound accumulated orbit error or certify periodicity.
Ordinary \code{mpf\_mul} can clip inputs differently, so exact split assembly
need not be bitwise identical to MT3. Zero or inputs shorter than two limbs
use ordinary products in the caller.

\paragraph{Worker ownership and handoff.}
Each square has high, low, and sum workers; its high worker assembles the
product after its two followers complete. The caller computes the cross
product's high part, while two workers compute its low and sum parts.
All operand snapshots, products, and assembly buffers are sized before workers
start and reused throughout the orbit. Each worker has a separate padded
atomic state with \code{Idle}, \code{Run}, and \code{Stop} values. Release
stores and acquire loads publish jobs and results, with processor pause
instructions while polling. Followers are dispatched before their lead,
preventing a lead from treating the previous Idle state as new completion.
There is at most one outstanding job per worker, so generation counters and
a shared completion counter are unnecessary. Pending arithmetic must finish
before snapshots are overwritten or workers are stopped.

\paragraph{Evaluated software pipeline.}
An experimental split interface started X-square work before sample publication
and derivative bookkeeping. The caller finished the current cross product and
formed $y_{n+1}$. After storing the completed $y_n^2$, the same Y snapshot
was prepared for $y_{n+1}^2$ while the caller consumed $x_n^2$ and performed
the full-precision real update. \Cref{fig:mt9-lookahead} shows the dependencies.
The next iteration started only the new X-square if its Y-square was already
pending. The prototype drained pending arithmetic before fallback, cancellation,
or destruction and suppressed lookahead at known termination and the final
iteration. Independent integer-oracle tests confirmed unchanged arithmetic.

Neither the early-dispatch experiment nor the full orbit-lookahead prototype
met the measurement criterion below. The split interface and experimental
schedule were removed. The retained MT9 \code{Compute} prepares both input
snapshots, dispatches all followers before the square leads, assembles the
cross product on the caller, and waits for both square leads before storing
the three outputs. Product assembly still overlaps other products within a
batch, but the retained orbit loop has no next-iteration square lookahead.

\begin{figure}[!htbp]
\centering
\begin{tikzpicture}[
  work/.style={draw, rounded corners, align=center, text width=2.7cm,
               minimum height=10mm, font=\small},
  dependency/.style={-{Latex[length=2mm]}, thick}]
  \node[work] (xsq) at (0,1.5) {$x_n^2$};
  \node[work] (cross) at (0,0) {$x_ny_n$};
  \node[work] (ysq) at (0,-1.5) {$y_n^2$};
  \node[work] (ynext) at (4.1,0) {Form $y_{n+1}$};
  \node[work] (consume) at (4.1,-2) {Consume $y_n^2$};
  \node[work, fill=blue!8] (lookahead) at (8.2,0) {Start $y_{n+1}^2$};
  \node[work, fill=blue!8] (real) at (8.2,-2) {Finish $x_{n+1}$};
  \draw[dependency] (cross) -- (ynext);
  \draw[dependency] (ynext) -- (lookahead);
  \draw[dependency] (ysq) -- (consume);
  \draw[dependency] (consume) -- (real);
  \draw[dependency] (consume.east) -- ++(0.45,0) |- (lookahead.south);
  \draw[dependency] (xsq.east) -- (9.85,1.5) |- (real.east);
  \node[align=center, font=\small] at (11.3,-0.7) {Shaded work\\may overlap};
\end{tikzpicture}
\caption[Imaginary-square lookahead dependencies]{Imaginary-square lookahead
  dependencies used by MT3 and evaluated in the MT9 prototype. The current cross product and
  Y-square must finish before Y's snapshot is replaced. Squaring the next Y
  can overlap the current real update; positions do not represent measured durations.}
\label{fig:mt9-lookahead}
\end{figure}

\paragraph{Measured pipeline evaluation.}
Measurements on October 11, 2026 used an AMD Ryzen 9 5950X with 16 physical
cores and 32 logical processors, Windows x64 Release, and 16,384 measured
iterations per fixture. The complete-orbit workload was $c=-0.1+0.1i$,
with benchmark-mode periodicity termination disabled. The product fixture used
fixed full-mantissa signed inputs and excluded construction and oracle checks
from timing; it is not a complete recurrence benchmark.

Each stage alternated baseline and candidate order across five pairs. The
retention rule required at least a 5\% median latency reduction in one MT9
fixture with no regression greater than 5\% in another. Results within two
percentage points of either threshold required ten additional pairs. Early
dispatch qualified for that extension, but its best gain across all fifteen
pairs was 4.23\% in the 16,384-bit product fixture; complete-orbit gains were
0.87\%, 1.85\%, and $-1.03$\% at increasing precision. The stage was rejected.

The orbit-lookahead prototype was compared against the original retained
schedule in a fresh set of five pairs. \Cref{tab:mt9-orbit-pipeline-timings}
reports complete-orbit medians and \cref{tab:mt9-product-pipeline-timings}
reports synchronous product medians. Positive gains indicate lower latency.
The unchanged MT3 control's median differences were below 1\% at every
precision. The timing runs included no concurrent build or CPU profiling.

\begin{table}[!htbp]
\centering
\caption[MT9 orbit-lookahead timings]{Complete-orbit MT9 timings, in microseconds
  per iteration, for five alternating baseline/prototype pairs.}
\label{tab:mt9-orbit-pipeline-timings}
\begin{tabular}{rrrr}
\hline
\textbf{Bits} & \textbf{Baseline ($\mu$s)} & \textbf{Lookahead ($\mu$s)} & \textbf{Gain (\%)} \\
\hline
16,384 & 9.288 & 9.450 & $-1.75$ \\
33,344 & 24.535 & 24.407 & 0.52 \\
100,096 & 98.362 & 99.002 & $-0.65$ \\
\hline
\end{tabular}
\end{table}

\begin{table}[!htbp]
\centering
\caption[MT9 product-control timings]{Synchronous three-product MT9 timings,
  in microseconds per batch, from the same five paired runs. The product
  fixture retains the original batch schedule in both binaries.}
\label{tab:mt9-product-pipeline-timings}
\begin{tabular}{rrrr}
\hline
\textbf{Bits} & \textbf{Baseline ($\mu$s)} & \textbf{Prototype ($\mu$s)} & \textbf{Gain (\%)} \\
\hline
16,384 & 7.424 & 7.601 & $-2.39$ \\
33,344 & 18.285 & 18.235 & 0.28 \\
100,096 & 72.416 & 73.082 & $-0.92$ \\
\hline
\end{tabular}
\end{table}

The proposed overlap does not remove the dependency of the next cross product
on the completed real component. Moving Y-square work earlier therefore need
not shorten the critical path when the current cross product dominates.
The measured result establishes no useful speedup for these workloads;
timing medians alone do not establish per-thread utilization or rule out gains
for other hardware or reference parameters.
